\subsection{Total categories, ordinary colimits, and the missing cofinality hypothesis}

\begin{result}{A.1}{Proved}{The two Grothendieck constructions and their comparison maps}
For a covariant diagram $D:A\to\mathrm{Cat}$, write
$\int_A^{\mathrm{opf}}D$ for the category with objects $(a,x)$,
$x\in D(a)$, and arrows
\[
 (a,x)\longrightarrow(b,y),\qquad
 (u:a\to b,\;v:D(u)x\to y).
\]
The superscript means the opfibrational construction. Its canonical functor
\[
 q_D:\int_A^{\mathrm{opf}}D\longrightarrow
                  \operatorname{colim}_{A}^{\mathrm{Cat}}D
\]
sends $(a,x)$ to its colimit image and $(u,v)$ to the image of $v$.
For a strict contravariant diagram $F:T^{\mathrm{op}}\to\mathrm{Cat}$,
write $\int_T^{\mathrm{fib}}F$ for the category whose arrows are
$(u:t\to s,v:x\to F(u)y)$. These are the conventions at VII ms.~99--110.
If $u:\mathbb N^{\mathrm{op}}\to T$ is a chosen sequence,
$D=F u^{\mathrm{op}}:\mathbb N\to\mathrm{Cat}$ and
$C_u=\int_{\mathbb N^{\mathrm{op}}}^{\mathrm{fib}}D$, there are specified
functors
\[
 C_u\xrightarrow{p}\int_T^{\mathrm{fib}}F=C,
 \qquad C_u\xrightarrow{q}\operatorname{colim}_{n}D(n)=C_\infty.
\]
There need not be a functor $C\to C_\infty$.
\end{result}
\begin{proof}
The colimit identifies the two images of $x$ and $D(u)x$, so the formula
for $q_D$ preserves identities and compositions. In the fibred case,
an arrow $(n,x)\to(m,y)$ has $n\geq m$ and an arrow
$x\to D(m\leq n)y$ in $D(n)$; its image is the corresponding arrow
$[x]\to[y]$ in the colimit. The functor $p$ is identity on fibre data
and applies $u$ to the base arrow. The formulas type every arrow in the
display. The manuscript correctly uses precisely this span at ms.~108;
inverting $q$ will give $C_\infty\to C$ in $\mathrm{Hot}$, rather than
an otherwise unspecified map in the reverse direction.
\end{proof}

\begin{result}{A.2}{Proved}{Filtered category colimits}
If $A$ is a small filtered category, then $q_D$ of A.1 belongs to
$W_\infty$ for every $D:A\to\mathrm{Cat}$, with no injectivity,
full-faithfulness, or fibrancy assumptions on its transition functors.
In particular $q:C_u\to C_\infty$ is in $W_\infty$ for every sequence.
The same universal assertion holds if each connected component of $A$
is filtered, and also for the empty category.
\end{result}
\begin{proof}
Fix an object $z$ of the category colimit. An object of $q_D/z$ is an
object $x\in D(a)$ together with an arrow $[x]\to z$ in the colimit.
Every finite list of objects, arrows, and equalities in a filtered
category colimit is represented at a common stage. For completeness,
construct a filtered category colimit by taking filtered colimits of
the object and arrow sets. Source, target, identities, and composition
descend because filtered colimits of sets commute with finite limits;
an equality of two represented arrows holds at a further stage.
The resulting category has the requisite universal property.

The category $q_D/z$ is nonempty by representing $z$ at some stage.
Given a finite diagram in $q_D/z$, choose a common stage $b$ representing
its objects, all its arrows to $z$, and all its commuting triangles.
Choose $z_b\in D(b)$ representing $z$. After passing to a further
stage if necessary, the arrows to $z_b$ give a cocone from that finite
diagram to $(b,z_b,\mathrm{id}_z)$ in $q_D/z$. Thus $q_D/z$ is filtered.
Its nerve is weakly contractible: the image of any finite simplicial
set in its nerve admits a cone after collecting its finite data, so
each sphere map becomes nullhomotopic. Quillen's Theorem A
\citep[Example 7.2.3.3, applied also to opposites]{ALurie2026} makes the
specified $q_D$ a nerve weak equivalence.

For the fibred sequence, take opposites. We have
\[
 C_u^{\mathrm{op}}=\int_{\mathbb N}^{\mathrm{opf}}D^{\mathrm{op}},
 \qquad
 (C_\infty)^{\mathrm{op}}
       =\operatorname{colim}_{n}D(n)^{\mathrm{op}}.
\]
Opposition preserves colimits and the weak homotopy type of a nerve.
The preceding argument therefore proves the second assertion. Both
constructions and their comparison split as coproducts over the
connected components of $A$. For $A=\varnothing$ both categories are
empty and the comparison is the identity.
\end{proof}

\begin{result}{A.3}{Proved}{An exact base condition uniform in the fibration}
Let $u:J\to T$ be a functor of small categories. The following
conditions are equivalent: every category $u/t$, $t\in T$, has weakly
contractible nerve; and, for every strict
$F:T^{\mathrm{op}}\to\mathrm{Cat}$, the projection
\[
 p_F:\int_J^{\mathrm{fib}}F u^{\mathrm{op}}
           \longrightarrow\int_T^{\mathrm{fib}}F
\]
belongs to $W_\infty$. In particular, for $J=\mathbb N^{\mathrm{op}}$,
this homotopy-initiality condition gives the specified isomorphism
$[p_F][q]^{-1}:C_\infty\to C$ in $\mathrm{Hot}$.
It is sufficient for one particular $F$ that $B(u/t)$ be weakly
contractible for every $t$ with $F(t)\ne\varnothing$.
\end{result}
\begin{proof}
For $c=(t,z)$ in $\int_T^{\mathrm{fib}}F$, the category $p_F/c$
maps to $u/t$ by forgetting its fibre data. Above $(j,\alpha:u(j)\to t)$
the remaining data are an object of
$F(u(j))/F(\alpha)z$. That category has terminal object
$\mathrm{id}_{F(\alpha)z}$. There is also an explicit section of the
forgetful functor, selecting this terminal object, and a natural
transformation from the identity to the section after the projection.
Explicitly, the given arrow $x\to F(\alpha)z$ supplies the
component from $(j,x,\alpha,x\to F(\alpha)z)$ to the selected object.
Thus $B(p_F/c)\simeq B(u/t)$ by actual functors and a natural
transformation. Quillen's Theorem A proves sufficiency.

For necessity, fix $t\in T$ and take the discrete-valued presheaf
$F_t(s)=\operatorname{Hom}_T(s,t)$. Its fibred construction is $T/t$,
with terminal object $(t,\mathrm{id}_t)$, and the pulled-back
construction is $u/t$. The universal hypothesis on $p_{F_t}$ forces
$B(u/t)$ to be weakly contractible. This proves the equivalence without
assuming that parallel reindexing functors are homotopic. Necessity is
claimed for the universal quantifier over all $F$, not for a particular
fibration or for the restricted subclass satisfying VII's homotopy
hypothesis.
\end{proof}

\begin{result}{A.4}{Refuted}{Preorder cofinality and parallel-map homotopies do not suffice}
The hypotheses printed at VII ms.~99--101 do not imply that
$C_u\to C$ or $C_\infty\to C$ is a weak equivalence. This remains false
when $T$ has one object and the fibre category is terminal. Thus
VIIb-C026 requires a condition such as A.3, or a proof specific to the
path fibration, beyond cofinality in the associated preorder.
\end{result}
\begin{proof}
Let $T=\mathbf B C_2$, the category with one object and endomorphisms
$1,s$, $s^2=1$.
Take $F$ constant at the terminal category and let
$u:\mathbb N^{\mathrm{op}}\to T$ send every arrow to $1$.
The preorder of $T$ is a singleton, the sequence is cofinal in it,
and all reindexing functors on the terminal fibre coincide.
Nevertheless $C=T$, $C_u=\mathbb N^{\mathrm{op}}$, and
$C_\infty=e$. The last two nerves are contractible whereas
$\pi_1(BT)=C_2$. The multiplication table and normalized chain
differentials are independently checked in
\texttt{code/a\_finite\_checks.py}; they give $H_1(BT;\mathbb Z)=\mathbb Z/2$.
Grothendieck already identifies the one-object obstruction himself at
VII ms.~109. It is not a correction to that page; it refutes the
insufficient earlier general hypothesis.
\end{proof}

\begin{result}{A.5}{Refuted}{The totally aspherical shape does not control colimits in Cat}
Let $A=\Delta_{\mathrm{inj}}$ have objects $[n]=\{0,\ldots,n\}$,
$n\geq0$, and strictly increasing functions. Let $D([n])$ be the
one-object category with endomorphism monoid $\mathbb N^{n+1}$.
An injection sends its standard generators to the indicated standard
generators. Then
\[
 \operatorname{colim}_{A}^{\mathrm{Cat}}D=\mathbf B\mathbb N,
 \qquad
 H_k\left(B\int_A^{\mathrm{opf}}D;\mathbb Z\right)=\mathbb Z
       \quad(k\geq0).
\]
Here $\mathbf B\mathbb N$ denotes the one-object category; its
classifying space has homotopy type $S^1$. Thus $q_D$ is not a weak
equivalence, despite the total $W_\infty$-asphericity of $A$.
\end{result}
\begin{proof}
The injections $[0]\to[n]$ identify all the standard generators in
any cocone, and connectedness identifies all the single objects.
A cocone into an arbitrary category is therefore exactly an object
with one chosen endomorphism. This is the universal property of the
one-object free monoid $\mathbb N$, proving the claimed category
colimit, without using a colimit in commutative monoids instead.

The homomorphism $\mathbb N^r\to\mathbb Z^r$ induces a weak
equivalence on classifying spaces. One direct proof is the covering
category with objects $\mathbb Z^r$ and a unique arrow $v\to w$ when
$w-v\in\mathbb N^r$: its underlying poset is filtered, its nerve is
contractible, and the free $\mathbb Z^r$-action has quotient nerve
$B(\mathbf B\mathbb N^r)$. Consequently the integral homology of $BD([n])$ is
$\Lambda^k\mathbb Z^{n+1}$. For $k\geq1$ the ordered wedge basis gives
an isomorphism of covariant $\mathbb Z A$-modules
\[
 H_k(BD(-);\mathbb Z)
   \cong\mathbb Z[\operatorname{Hom}_{A}([k-1],-)].
\]
These are projective representable modules, whose colimit is
$\mathbb Z$ and whose higher derived colimits vanish. For $k=0$ the
module is constant $\mathbb Z$. The functor $S([n])=[n+1]$ which
adds a new last element admits natural transformations
$\mathrm{id}_A\Rightarrow S\Leftarrow\operatorname{const}_{[0]}$,
using respectively the original elements and the new last element.
Therefore $BA$ is contractible and the higher derived colimits of
this degree-zero module also vanish.

Thomason's homotopy-colimit comparison, explicitly restated in
\citet[Section~3.1, equation~(9)]{ACegarra2020}, identifies
the total nerve with the homotopy colimit. The bar double complex
gives the first-quadrant spectral sequence
\[
 E^2_{p,q}=L_p\operatorname{colim}_{A}H_q(BD(-);\mathbb Z)
 \ \Longrightarrow\ H_{p+q}\left(B\int_A^{\mathrm{opf}}D;\mathbb Z\right).
\]
Every column except $p=0$ vanishes and $E^2_{0,q}=\mathbb Z$ for every
$q\geq0$, so it has no differential or extension ambiguity. In
particular the source of $Bq_D$ has $H_2=\mathbb Z$ and its target
has $H_2=0$. This verifies the counterexample proposed after TOTASPH.
It does not require the full exactness theorem for arbitrary
abelian-group-valued diagrams. The accompanying program checks the
ordered wedge/representable identification and the colimit relations
through $[5]$; the proof above, rather than extrapolation from those
finite checks, establishes the assertion in all degrees.

We verify the total-asphericity premise in A.6 below, recording the
argument rather than treating the earlier AI-produced answer as an
authority.
\end{proof}

\begin{result}{A.6}{Proved}{Verification of the Delta-injective premise}
The category $\Delta_{\mathrm{inj}}$ has exact colimits of
abelian-group-valued diagrams and is totally $W_\infty$-aspherical,
but is not filtered. These are the assertions of TOTASPH C.7--C.10
used in A.5.
\end{result}
\begin{proof}
We first check the finite algebraic construction which underlies the
inherited proof. For every $n\geq0$ there is an integral combination
$t$ of injections $[n]\to[M]$, of augmentation one, such that $th$
is independent of the injection $h:[m]\to[n]$ for each fixed $m$.
Here negative coefficients are allowed. To construct it, for each
$k\geq1$ choose a homogeneous integral polynomial
$C_{n,k}(z_0,\ldots,z_{n+1})$ having value one at all variables equal
to one and such that $y_i^kD_iC_{n,k}$ is independent of $i$.
The map $D_i$ identifies $z_i,z_{i+1}$ and relabels the remaining
variables $y_0,\ldots,y_n$. Set $C_{0,k}=1$. Inductively put
$Q=C_{n-1,2k}$ and prescribe
\[
 D_i C_{n,k}=Q\prod_{l\ne i}y_l^k.
\]
These prescriptions agree on pairwise intersections. For adjacent
indices both products agree after the second identification; for
nonadjacent indices the required identity is precisely
$w_p^{2k}D_pQ=w_q^{2k}D_qQ$. To see that agreement is sufficient
over the integers, change coordinates to
$x_i=z_i-z_{i+1}$ and $u=z_{n+1}$. For compatible prescribed residues
$F_i$ modulo $(x_i)$, let $\rho_S$ set $x_i=0$ for $i\in S$ and set
\[
 C=\sum_{\varnothing\ne S\subseteq\{0,\ldots,n\}}
             (-1)^{|S|+1}\rho_S F_i\qquad(i\in S).
\]
Compatibility makes each term independent of $i\in S$; cancellation
modulo $(x_j)$ leaves $F_j$. Homogeneity is preserved. The degree
satisfies $d(n,k)=d(n-1,2k)+nk=k(2^{n+1}-n-2)$.

An injection $g:[n]\to[M]$ corresponds to the monomial with exponents
$g(0),g(1)-g(0)-1,\ldots,g(n)-g(n-1)-1,M-g(n)$, of total degree
$M-n$. Precomposition with the coface omitting $i$ acts on this
polynomial by $P\mapsto y_iD_iP$. Choose $P=C_{n,1}$ and
$M=n+d(n,1)$. All $t\delta^i$ are equal. The coface identity
$\delta^{b+1}\delta^a=\delta^a\delta^b$ then shows that this common
value equalizes all injections into $[n-1]$, by downward induction;
factoring a proper injection through a coface proves the assertion
for every $h$. The augmentation is the evaluation at all variables
equal to one, hence is one.

Let $\operatorname{Aff}_{\mathbb Z}(A)$ have the objects of $A$ and
augmentation-one finite integral arrow combinations as its arrows.
Any two objects of $A$ have a common successor. The combination $t$
just constructed equalizes any two parallel affine arrows, because
both have coefficient sum one. Thus $\operatorname{Aff}_{\mathbb Z}(A)$
is filtered. An abelian-group-valued functor on $A$ extends uniquely
by linearity to this affine category. Cocones over the two diagrams
are the same: respecting ordinary arrows implies respecting affine
combinations of augmentation one. The colimits are therefore
canonically the same. Since filtered colimits of abelian groups are
exact, so are colimits over $A$.

For a covariant $H:A\to\mathrm{Set}$ and a diagram $V$ of abelian
groups on its category of elements, the assignment
$(\Sigma_H V)(a)=\bigoplus_{x\in H(a)}V(a,x)$ is exact and has
$\operatorname{colim}_A\Sigma_H V
 =\operatorname{colim}_{\operatorname{El}_A H}V$ by the cocone
universal property. Hence colimits on every connected component $E$
of every such element category are exact, and the bar resolution
gives $H_j(BE;\mathbb Z)=0$ for $j>0$. The same argument applies to
categories of elements over $E$: their iterated element construction
is an element construction over $A$. For a cyclic subgroup
$L\leq\pi_1(BE)$, the associated covering is the element category
of a locally constant set-valued functor on $E$. It has
$H_1=L_{\mathrm{ab}}=L$ and must be acyclic. All cyclic subgroups
are thus trivial. Each $BE$ is simply connected and integrally
acyclic, hence weakly contractible by Hurewicz.

If $f:A\to B$ is ordinary final, its comma category $b\backslash A$
is a nonempty connected element category for every $b$, and therefore
is weakly contractible. Quillen's Theorem A makes $Bf$ a weak
equivalence. We already proved $BA$ contractible, so $BB$ is
contractible. This is total asphericity in the final-image convention
of the prompt. Finally, the two arrows $[0]\rightrightarrows[1]$
cannot be equalized by postcomposition with an injection, proving
nonfilteredness.
\end{proof}

\begin{result}{A.7}{Proved}{Bounds on the class of universally correct Cat-colimit shapes}
Let $\mathcal C$ be the class of nonempty connected small categories
$A$ for which $q_D\in W_\infty$ for every $D:A\to\mathrm{Cat}$.
Then
\[
 \{\text{filtered categories}\}\ \subseteq\ \mathcal C
 \ \subsetneq\ \{\text{totally }W_\infty\text{-aspherical categories}\}.
\]
The right inclusion is strict by A.5--A.6. For arbitrary, possibly
disconnected $A$, the comparison property is componentwise and the
empty shape satisfies it.
\end{result}
\begin{proof}
The first inclusion is A.2. For the second, apply the comparison
property to every discrete-valued diagram $H:A\to\mathrm{Set}$.
Its Grothendieck construction is $\operatorname{El}_A H$ and its
colimit is the discrete set $\operatorname{colim}_A H$. Every
connected component of the element category must therefore be
weakly contractible. In particular the constant singleton diagram
makes $BA$ contractible. For an ordinary final $f:A\to B$, use
$H_b(a)=\operatorname{Hom}_B(b,f(a))$ and Quillen's Theorem A exactly
as at the end of A.6. Then $BB$ is contractible, which proves total
asphericity. The componentwise and empty assertions follow from
the coproduct decompositions in A.2.
\end{proof}

\begin{result}{A.8}{Open}{Exact classification of Cat-colimit shapes}
We do not determine whether every member of $\mathcal C$ in A.7 is
filtered. More generally, an intrinsic criterion for membership in
$\mathcal C$, sharper than the proven bounds, is unresolved in this
paper. The finite-object and weakly-terminal scopes reduce to
filteredness by the independently checkable total-asphericity
criterion TOTASPH C.13; we do not need that additional criterion in
the preceding arguments. The first scope not decided by those
results consists of infinite categories with no weakly terminal
object. No smallest counterexample is asserted.
\end{result}
The simplicial-set criterion cannot prove the missing implication:
the map it controls has target $\operatorname{colim}_A ND$, whereas
the target needed here is $N\operatorname{colim}^{\mathrm{Cat}}_A D$.
A.5 explicitly separates these targets. Nor does exactness of
abelian-group colimits suffice, since A.6 verifies that property
for the same counterexample. The searches and the exact status of
the sources are recorded in the report.

\begin{result}{A.9}{Proved}{A finite obstruction to replacing Hot by Cat hot}
Let $P$ be the face poset of the boundary of a triangle, with three
vertices and three edges, and let $\operatorname{sd}P$ be the poset
of nonempty chains in $P$, ordered by inclusion. The last-element
map $\ell:\operatorname{sd}P\to P$ is a weak equivalence and is not
invertible in $\mathrm{Cat\,hot}$. There are exactly $456$ monotone
maps $P\to\operatorname{sd}P$, and all induce degree zero on the
circle.
\end{result}
\begin{proof}
The objects of $P$ are the six faces
\[
 \{0\},\quad\{1\},\quad\{2\},\quad
 \{0,1\},\quad\{1,2\},\quad\{0,2\}.
\]
Its arrows are their inclusions. The objects of $\operatorname{sd}P$
are their nonempty chains, with inclusions as arrows. Compositions
are the unique arrows given by inclusion.
Their nerves are respectively cycles with six and twelve edges.
The standard subdivision homeomorphism and the simplexwise linear
homotopy to the last vertex show that $B\ell$ has degree one.
Every monotone map in the reverse direction gives a simplicial map
which sends each of the six source edges to one target edge or a
vertex. The total winding has absolute value at most $6/12<1$;
being integral, its degree is zero. It therefore cannot be a homotopy
inverse even after passing to nerves. Each natural transformation
induces a homotopy of nerves, and the same holds for any finite
zigzag, so no inverse exists in $\mathrm{Cat\,hot}$.
The program \texttt{code/a\_finite\_checks.py} verifies both partial
orders, enumerates all vertex images and all compatible edge images,
and computes every winding number. Its output is $456$ maps of
degree zero and degree one for $\ell$. This verifies POLES D.6 and
provides a finite warning for VIIc-A046 and V-I066; it does not by
itself decide the particular path comparison $q_L$.
\end{proof}
